\section{Related Work}
\label{sec:related}

Early VQA systems classified over a fixed set of answers, whereas current generative VLMs produce free-form answers with an autoregressive language
model~\cite{li2023blip2,alayrac2022flamingo,liu2023visualinstructiontuning}. Vietnamese datasets include ViVQA~\cite{tran2021vivqa}, OpenViVQA~\cite{nguyen2023openvivqa},
ViCLEVR~\cite{tran2023viclevr}, ViOCRVQA~\cite{pham2024viocrvqa} and ViTextVQA~\cite{nguyen2025vitextvqa}. AutoViVQA~\cite{autovivqa2026} contributes 37{,}077
questions on 19{,}411 COCO images~\cite{lin2015coco}, generated by large language models,
filtered by an ensemble of validators and annotated with reasoning categories, with an
emphasis on relational, causal and multi-step reasoning. On the modelling side,
BARTPhoBEiT~\cite{tran2023bartphobeitpretrainedsequencetosequenceimage} and
AViVQA-TranConI~\cite{nguyen2024tranconi} combine pre-trained visual and Vietnamese
language representations, Vintern-1B~\cite{doan2024vintern1b} brings an instruction-tuned
VLM to Vietnamese, and ViMoE-VQA~\cite{vimoevqa} routes multimodal tokens to sparse experts
for generative VQA on AutoViVQA. These approaches either train a new model or fine-tune an
existing VLM on multi-tile inputs; in contrast, we keep every pre-trained weight of
Vintern-1B frozen and use a single image view.

Keeping the vision encoder and the language model frozen and training only a connector is a
common way to build VLMs efficiently, and connectors differ mainly in what they must learn.
Frozen~\cite{tsimpoukelli2021frozen} learns a visual prefix from a global image
representation. The Q-Former of BLIP-2~\cite{li2023blip2} learns a set of queries that
compress the patch features, which preserves details but requires a large module (188M
parameters) pre-trained on millions of image--text pairs. LLaVA~\cite{liu2023visualinstructiontuning,liu2024llava15}
projects every patch token with a small MLP, which preserves locality but passes hundreds of
tokens to the decoder; Honeybee~\cite{cha2024honeybee} proposes locality-preserving
abstractors, and InternVL~\cite{chen2024internvl}, on which Vintern-1B is built, merges
neighbouring patches before its MLP projector. In all of these models, the mapping from
patches to the decoder is learned during large-scale pre-training. ViBridge-VQA differs in two
respects: it does not learn this local mapping again but reuses the projector already
pre-trained in the VLM, and it trains only a small global bridge (12.9M parameters) on about
26k task questions. 

The number of visual tokens largely determines the cost of the decoder. Inference-optimal
scaling of VLMs~\cite{li2024inferenceoptimal} shows that a few tokens can suffice for many
tasks, and Matryoshka multimodal models~\cite{cai2024matryoshka} train a model to work at
several nested pooled token budgets and show that the required budget depends on the task.
These approaches train or re-train the model for each budget. In contrast, we keep the model
frozen and obtain each budget simply by pooling the output of the frozen projector, and we
choose among budgets with explicit criteria, namely the knee of the Pareto
front~\cite{satopaa2011kneedle} and TOPSIS~\cite{hwang1981topsis}. Parameter-efficient
methods such as adapters~\cite{houlsby2019adapters}, prefix tuning~\cite{li2021prefix} and
LoRA~\cite{hu2022lora} reduce the cost of adaptation from the other side, by training a small
number of additional weights in the decoder. We use decoder LoRA both as a baseline and as a
complement, and ask whether a small training budget is better spent on the visual tokens or on
adapting the decoder.
